\chapter{What a Test Schedule is}\label{ch:what-a-schedule-is}

ATE runs a scripted test against a unit. The script is a \textbf{Test Schedule}: a Lua file
that describes, in order, the \textbf{Test Steps} to run. ATE reads the file, executes the
steps one after another, records every result in a database and writes a \textbf{Test Report}.
A schedule never drives ATE; ATE drives the schedule.

The vocabulary is fixed. A \textbf{UUT} is the unit under test, identified by its UUT Type and
serial number. The \textbf{Operator} is the person running the Test Run. A \textbf{Station} is
one installation of ATE on one computer. A \textbf{Device} is a connection ATE talks through:
a TCP socket, a serial connection, a UDP datagram channel, or a VISA resource.
\textbf{Equipment} is a record of what was on the bench, declared by the schedule so it appears
in the Test Report. A \textbf{Media Release} is a set of Lua files delivered together, described
by a \textbf{Test Definition}. The names are the ones in GLOSSARY.md: the handle a step gets from
\texttt{ATE.open} is a Device, never an instrument or a port, and \emph{instrument} means the
physical equipment it reaches.

\section{How much Lua this manual assumes}

This manual assumes you can read basic Lua: \texttt{local}, functions, tables, \texttt{if},
\texttt{for}, \texttt{while}, string concatenation with \texttt{..}, and \texttt{nil}. ATE's
schedule syntax has three idioms worth naming:

\begin{itemize}
  \item \texttt{a:foo():bar()} chains method calls, because every setter returns the object it
    was given.
  \item A params table is a Lua table: \texttt{\{ host = "127.0.0.1", port = 5025 \}}. Keys are
    written \texttt{key = value}, and a table with no value for a key reads as \texttt{nil}.
  \item \texttt{function() \dots end} is an anonymous function; a step's action is one, and a
    Sub-schedule's parameters are passed to it as \texttt{local kv = \dots}.
\end{itemize}

Nothing else in the language is used by the examples in this manual.

ATE targets Lua 5.5. The runtime is a normal Lua 5.5 interpreter with a few entry points
removed (chapter~\ref{ch:sandbox}), so plain Lua knowledge transfers. A step's action runs
synchronously on the worker thread. Coroutines are not taken away, and using one is not
forbidden, but the count hook does not reach a coroutine: a busy loop inside one is one of the
two documented exceptions to the ten-second abort rule (chapter~\ref{ch:errors}).

\section{The shape of a Test Schedule}

A Test Schedule is a Lua file with the extension \texttt{.tslua}. Its top level runs when the
file is loaded, before any Test Run exists, and it must \emph{return} a schedule object built
with \texttt{ScheduleParams}. Nothing in the file executes a Test Step itself. Two things make a
file fail to load, and the manual calls both a load error: an \texttt{ATE} call that needs a run
at the top level, which raises with the function's name (chapter~\ref{ch:sandbox}), and a
schedule that fails validation, such as a duplicate Step Number or a step with no action.

\section{Where a schedule lives}

A schedule is shipped inside a Media Release: a directory of Lua files plus one
\texttt{.atedef} file that lists every file in it with its checksum (chapter~\ref{ch:modules}).
ATE finds Test Definitions by searching the directories named in the application config,
loads every visible schedule, and shows their names and descriptions on the Run screen.
The Operator picks a schedule, types the UUT Type and serial number, and presses Start.

\cando{name the pieces of a Test Schedule and say where each one comes from: the
\texttt{.tslua} file, the \texttt{.atedef} that ships it, and the Station's config that
makes it available to run.}
